\section{Preliminaries}
\subsection[$q$-calculus]{$\boldsymbol{q}$-calculus}

We give a short introduction to $q$-derivatives, $q$-integration and $q$-special functions, see \cite{MR1052153, JACKSON, Koekoek, MR1448685}. The report \cite{Koekoek} that will be referred to often is also included in the book \cite{MR2656096}. For $u$ a~number, and $q$ the deformation parameter, $0<q<1$, we def\/ine the $q$-deformation of $u$ as
\begin{gather*}
[u]_q=\frac{q^u-1}{q-1}.
\end{gather*}
It is clear that $\lim\limits_{q\to 1}[u]_q=u$. We also def\/ine
\begin{gather*}
(u;q)_k=(1-u)(1-qu)\cdots\big(1-q^{k-1}u\big)\qquad \mbox{and}\qquad (u;q)_\infty=\prod_{k=0}^\infty \big(1-uq^k\big).
\end{gather*}


The $q$-derivative of a function $f(t)$ is def\/ined by
\begin{gather}
\label{qder}
\partial^q_t (f(t)) = \frac{f(qt)-f(t)}{(q-1)t}.
\end{gather}
This operator satisf\/ies the generalized Leibniz rule
\begin{gather}
\label{Leibniz}
%\label{Leibniz2}
\partial_t^q(f_1(t)f_2(t)) = \partial_t^q(f_1(t))f_2(t)+f_1(qt)\partial_t^q(f_2(t)).
\end{gather}


The $q$-integration on an interval $[0,a]$ with $a\in \mR$ is given by
\begin{gather*}
%\label{1Dqint}
\int_{0}^af(t)\,d_qt = (1-q)a\sum_{k=0}^\infty f\big(q^ka\big)q^k.
\end{gather*}
The inf\/inite $q$-integral can be def\/ined in several ways, determined by a parameter $\gamma\in\mR\backslash \{0\}$,
\begin{gather}
\label{gammaint}
\int_0^{\gamma\cdot\infty}f(t)d_qt = (1-q)\gamma\sum_{k=-\infty}^\infty f\big(q^k\gamma\big)q^k=\lim_{l\to+\infty}\int_0^{q^{-l}\gamma}d_qt f(t).
\end{gather}
The positive $(\gamma\in\mR^+)$ or negative $(\gamma\in\mR^-)$ inf\/inite integral is a function of $\gamma$, however from the def\/inition it is clear that $\int_0^{\gamma\cdot\infty}=\int_0^{q\gamma\cdot \infty}$. So $\partial_\gamma^q\int_0^{\gamma\cdot\infty}d_qt=0$ holds which means the integral is a $q$-constant. The integral is the inverse of dif\/ferentiation,
\begin{gather}
\label{partint1}
\int_0^{a}\left[\partial_t^qf(t)\right]d_qt = f(a)-f(0).
\end{gather}




The $q$-factorial of an integer $k$ is given by $[k]_q!=[k]_q[k-1]_q\cdots [1]_q$ and satisf\/ies $[k]_q!=(q;q)_k/(1-q)^k$. This can be generalized to the $q$-Gamma function $\Gamma_q(t)$ for $t>0$ satisfying $\Gamma_q(t+1)=[t]_q\Gamma_q(t)$, see e.g.~\cite[formula (1.10.1)]{MR1052153}. The $q$-exponentials are def\/ined as
\begin{gather*}
%\label{exp1}
e_q(t)=\sum_{j=0}^\infty\frac{t^j}{[j]_q!} \qquad \mbox{and} \qquad E_q(t)=e_{q^{-1}}(t)=\sum_{j=0}^\infty q^{\frac{1}{2}j(j-1)}\frac{t^j}{[j]_q!}.
\end{gather*}
Note that a dif\/ferent notation for the exponentials is used compared to \cite{MR2595267}. The relation $e_q(t)E_q(-t)=1$ holds and the derivatives are given by
\begin{gather}
\label{aflexp}
\partial_t^qe_q(t)=e_q(t) \qquad \mbox{and}\qquad \partial_t^qE_q(t)=E_q(qt).
\end{gather}
For $q<1$ the series $E_q(t)$ converges absolutely and uniformly everywhere and $e_q(t)$ in the area $|t|<\frac{1}{1-q}$. The function $e_q(t)$ can be analytically continued to $\mC\backslash \{\frac{q^{-k}}{1-q}\}$ as $1/(E_q(-t))$. The zeroes of the $q$-exponential $E_q$ are
\begin{gather}
\label{zeroes}
E_{q}\left(-\frac{q^{-k}}{1-q}\right) = 0\qquad \mbox{for} \ \ k\in\mN.
\end{gather}
This follows from the inf\/inite product representation $E_q(\frac{t}{1-q})=(-t;q)_\infty$, see \cite[formula (1.3.16)]{MR1052153}. This implies that the relation
\begin{gather}
\label{intinfeind}
\int_0^{\frac{1}{\sqrt{1-q}}\cdot\infty}d_qt f(t) E_{q^2}\left(\frac{-t^2}{1+q}\right)=\int_0^{\frac{1}{\sqrt{1-q}}}d_qt f(t) E_{q^2}\left(\frac{-t^2}{1+q}\right)
\end{gather}
holds.



The $q$-Hermite polynomials are given by
\begin{gather}
\label{Hermite}
H_k^q(t) = \sum_{j=0}^{\lfloor k/2\rfloor}(-1)^j\frac{[k]_q!}{[k-2j]_q![j]_{q^2}!}q^{j(j+1)}((q+1)t)^{k-2j}
\end{gather}
and related to the discrete $q$-Hermite~I polynomials $h_k(x;q)$ in \cite[Section~3.28]{Koekoek}  by
\begin{gather*}
H_k^q(t) = q^k\left(\frac{1+q}{1-q}\right)^{k/2}h_k\Big(q^{-1}\sqrt{1-q^2}t;q\Big).
\end{gather*}


We introduce the $q$-Laguerre polynomials for $\alpha>-1$ in the normalization of \cite[p.~24]{MR2595267},
\begin{gather}
\label{qLaginv}
\cL^{(\alpha)}_j\big(u|q^{-2}\big) = q^{-j(j+1+2\alpha)}\sum_{i=0}^jq^{2i(i+\alpha)}\frac{(-u)^i}{[j-i]_{q^2}![i]_{q^2}!}\frac{(q^{2i+2\alpha+2};q^2)_{(j-i)}}{(1-q^2)^{j-i}}.
\end{gather}
They can also be def\/ined using the $q$-Gamma function since $\frac{(q^{2i+2\alpha+2};q^2)_{(j-i)}}{(1-q^2)^{j-i}}=\frac{\Gamma_{q^2}(j+\alpha+1)}{\Gamma_{q^2}(i+\alpha+1)}$ holds. They are connected with the $q$-Laguerre polynomials $L_j^{(\alpha)}(u;q)$ from \cite[Section~3.21]{Koekoek}  by
\begin{gather*}
\cL_j^{(\alpha)}\big(u|q^{-2}\big) = q^{-j(j+1+2\alpha)}L_j^{(\alpha)}\big(\big(1-q^2\big)u;q^2\big)
\end{gather*}
and to the $q$-Laguerre polynomials in \cite{MR0618759} by the same formula with a substitution  $(1-q^2)u\to u$ in the right hand side of the formula.

The substitution $q\to q^{-1}$ yields,
\begin{gather}
\label{qLag}
\cL^{(\alpha)}_j\big(u|q^2\big) = \sum_{i=0}^jq^{(j-i)(j-i+1)}\frac{(-u)^i}{[j-i]_{q^2}![i]_{q^2}!}\frac{(q^{2i+2\alpha+2};q^2)_{(j-i)}}{(1-q^2)^{j-i}}.
\end{gather}
These polynomials are related to the little $q$-Laguerre polynomials (the Wall polynomials) $p_j(u;a|q)$, see \cite[Section~3.20]{Koekoek}, by
\begin{gather*}
\cL^{(\alpha)}_j\big(u|q^2\big) = q^{j(j+1)}\frac{(q^{2\alpha+2};q^2)_j}{(1-q^2)^j[j]_{q^2}!}p_j\left(\big(q^{-2}-1\big)u;q^{2\alpha}|q^2\right).
\end{gather*}
The $q$-Laguerre polynomials in equation \eqref{qLag} satisfy the orthogonality relation
\begin{gather*}
\int_0^{\frac{1}{1-q^2}}d_{q^2}u\,u^\alpha\cL_j^{(\alpha)}\big(u|q^2\big)\cL_k^{(\alpha)}\big(u|q^2\big)E_{q^2}(-u)=
\delta_{jk}q^{2(j+1)(\alpha+1+j)}\frac{\Gamma_{q^2}(j+\alpha+1)}{[j]_{q^2}!},
\end{gather*}
see \cite[equation (3.20.2)]{Koekoek}. Using the calculation rules in~\cite[Lemma 10]{MR2595267}, this can be rewritten~as
\begin{gather}
 \int_0^{\frac{1}{\sqrt{1-q}}}d_qr\, r^{2\alpha+1}\cL_j^{(\alpha)}\left(\frac{r^2}{1+q}|q^2\right)\cL_k^{(\alpha)}\left(\frac{r^2}{1+q}|q^2\right) E_{q^2}\left(\frac{-r^2}{1+q}\right)\nonumber\\
\qquad{} =\delta_{jk}\frac{q^{2(j+1)(j+\alpha+1)}\Gamma_{q^2}(j+\alpha+1)(1+q)^{\alpha}}{[j]_{q^2}!}.\label{orthLag1}
\end{gather}
One of the orthogonality relations for $q$-Laguerre polynomials in equation~\eqref{qLaginv} is
\begin{gather}
 \int_0^{\gamma\cdot\infty}d_qt\, t^{2\alpha+1}\cL_j^{(\alpha)}\left(\frac{q^2t^2}{1+q}|q^{-2}\right)\cL_k^{(\alpha)}\left(\frac{q^2t^2}{1+q}|q^{-2}\right)
 e_{q^2}\left(\frac{-q^2t^2}{1+q}\right)\nonumber\\
\qquad{} =\delta_{jk}\frac{q^{-(j+\alpha)(j+\alpha+1)}\Gamma_{q^2}(j+\alpha+1)
(1+q)^{\alpha}}{[j]_{q^2}!\,q^{(j+1)(j+2\alpha+2)}}d\left(\frac{\gamma}{\sqrt{1+q}},\alpha\right),\label{orthLag2}
\end{gather}
see \cite[equation~(3.21.3)]{Koekoek}, with
\begin{gather*}
d\left(\frac{\gamma}{\sqrt{1+q}},\alpha\right) = \frac{q^{(\alpha+1)(\alpha+2)}}{(1+q)^\alpha\Gamma_{q^2}(\alpha+1)}\int_0^{\gamma\cdot\infty}d_qt\, t^{2\alpha+1}\,e_{q^2}\left(\frac{-q^2t^2}{1+q}\right).
\end{gather*}
The function $d(\lambda,\alpha)$ therefore satisf\/ies
\begin{gather*}
d\big(\sqrt{\gamma},\alpha\big) = \frac{q^{\alpha(\alpha+1)}}{\Gamma_{q^2}(\alpha+1)}\int_0^{\gamma\cdot\infty}d_{q^2}u\,u^\alpha e_{q^2}(-u).
\end{gather*}
Partial integration implies that this function satisf\/ies $d(\lambda,\alpha+1)=d(\lambda,\alpha)$ for $\alpha>-1$. The explicit expression for $d$ can be found from \cite[equation~(3.21.3)]{Koekoek}.

\begin{remark}
The $q$-Laguerre polynomials do not form a complete orthogonal system for the Hilbert space corresponding to the measure in equation \eqref{orthLag2}. In \cite{MR1803885} the compliment of the basis is constructed. The corresponding functions derived in~\cite[Section~4]{MR1803885} (with a suitable renormalization) will therefore be annihilated by the $q$-Hankel transform $\cH_\nu^{q,\gamma}$ in the subsequent Def\/inition~\ref{defHankel}. This follows from the same calculation that leads to the subsequent equation~\eqref{preHank2}.
\end{remark}

The $q$-Gegenbauer polynomials, see \cite[equation~(2.19)]{MR1340930}, are given by
\begin{gather}
\label{qGeg}
C^\lambda_n(q;t)=\sum_{j=0}^{\lfloor\frac{n}{2}\rfloor}\frac{(-1)^jq^{j(j-1)}}{[j]_{q^2}![n-2j]_q!}
\frac{(q^{2\lambda};q^2)_{n-j}}{(1-q^2)^{n-j}}((1+q)t)^{n-2j}.
\end{gather}
They are big $q$-Jacobi polynomials on $[-1,1]$ with the two parameters equal to $\lambda-\frac{1}{2}$, see~\cite[equation~(2.26)]{MR1340930}.

For $\nu>-1$, the f\/irst and second $q$-Bessel function, introduced by Jackson, see~\cite{MR0649849, JACKSON}, are given by
\begin{gather}
\label{Bessel1}
J_\nu^{(1)}\big(x|q^2\big)=\left(\frac{x}{1+q}\right)^\nu\sum_{i=0}^\infty\frac{(-1)^i}{[i]_{q^2}!
\Gamma_{q^2}(i+\nu+1)}\left(\frac{x}{1+q}\right)^{2i} \qquad \mbox{for} \ \ |x|<\frac{1}{1-q}
\end{gather}
and
\begin{gather}
\label{Bessel2}
J_\nu^{(2)}\big(x|q^2\big)=q^{\nu^2}\left(\frac{x}{1+q}\right)^\nu\sum_{i=0}^\infty\frac{q^{2i(i+\nu)}(-1)^i}{[i]_{q^2}!
\Gamma_{q^2}(i+\nu+1)}\left(\frac{x}{1+q}\right)^{2i}.
\end{gather}
$J_\nu^{(1)}(x|q^2)$ is analytical in the area $|x|<\frac{1}{1-q}$ and $J_\nu^{(2)}(x|q^2)$ is analytical on $\mR^+$.
The f\/irst $q$-Bessel function can be analytically continued by the relation
\begin{gather*}
J_\nu^{(1)}\big(x|q^2\big) = \frac{1}{q^{\nu^2}}e_{q^2}\left(-\frac{1-q}{1+q}x^2\right)J_{\nu}^{(2)}\big(x|q^2\big),
\end{gather*}
see \cite[Exercise~1.24]{MR1052153}, which is def\/ined for all~$x$. Since $x^{-\nu}J_\nu^{(2)}(x|q^2)$ is an entire function the formula above implies that $x^{-\nu}J_\nu^{(1)}(x|q^2)$ is analytic on $\mC$ outside the poles $\{\pm i q^{-k}(1-q)^{-1}|k\in\mN\}$. Therefore $x^{-\nu}J_\nu^{(1)}(x|q^2)$ is well-def\/ined and analytic for~$x\in\mR$.



These $q$-Bessel functions are related to the $J^{(i)}_\nu(x;q)$ in \cite[(1.13) and (1.17)]{MR0649849}  or \cite[Exercise~1.24]{MR1052153} by
\begin{gather*}
J^{(1)}_\nu\big(x|q^2\big)=J^{(1)}_\nu\big({2(1-q)}x;q^2\big),\qquad J^{(1)}_\nu\big(x|q^2\big)=q^{\nu^2}J^{(1)}_\nu\big(2(1-q)x;q^2\big).
\end{gather*}



The generating functions for the $q$-Laguerre polynomials are given by
\begin{gather}
\label{BesselLag1}
J_\alpha^{(1)}\big(rt|q^2\big) = \left(\frac{rt}{1+q}\right)^{\alpha}\sum_{j=0}^\infty\frac{\cL_j^{(\alpha)}\big(\frac{r^2}{1+q}|q^2\big)}{\Gamma_{q^2}(\alpha+j+1)}
\frac{t^{2j}}{(1+q)^j}e_{q^2}\left(-\frac{q^2t^2}{1+q}\right),\\
\label{BesselLag2}
J_\alpha^{(2)}\big(qrt|q^2\big) = \left(\frac{rt}{1+q}\right)^{\alpha}\sum_{j=0}^\infty\frac{q^{(j+\alpha)(j+1+\alpha)}\cL_j^{(\alpha)}
\big(\frac{q^2t^2}{1+q}|q^{-2}\big)}{\Gamma_{q^2}(\alpha+j+1)}\frac{r^{2j}}{(1+q)^j}E_{q^2}\left(-\frac{r^2}{1+q}\right).
\end{gather}
This follows from direct calculations, they are equivalent to~\cite[formulas~(3.20.11) and (3.21.13)]{Koekoek}.



\subsection[The Howe dual pair and harmonic oscillator on $\mR^m_q$]{The Howe dual pair and harmonic oscillator on $\boldsymbol{\mR^m_q}$}

Quantum spaces are spaces where the variables have braid statistics. The commutation relations are generalizations of the bosonic or fermionic ones by an $\hat{R}$-matrix. The algebra of functions on a quantum space can be seen as the algebra $\cO$ of formal power series in non-commuting variables $x^1,\dots,x^m$,
\begin{gather*}
\cO = \mC[[x^1,\dots,x^m]]/I
\end{gather*}
with $I$ the ideal generated by the commutation relations of the variables. We consider quantum spaces which satisfy the Poincar\'e--Birkhof\/f--Witt property, which states that the dimension of the space of homogeneous polynomials of a certain degree is the same as in the commutative case. Superspaces, for instance, do not satisfy this property.

We focus on the case of the quantum Euclidean space $\mR^m_q$. The relations for the variables can e.g.\ be found in \cite{MR1328733, MR2006509, MR1361741}. We denote by $\cO_q$ the algebra of formal power series for the specif\/ic case of the quantum Euclidean space. The quantum Euclidean space can be def\/ined by the $\hat{R}$-matrix of the quantum orthogonal group $O_q(m)$, see~\cite{MR1097174, MR1239953}. The matrix $\hat{R}\in\mC^{(m\times m)\times(m\times m)}$ can be expressed in terms of its projection operators as $\hat{R}=q\cP_S-q^{-1}\cP_A+q^{1-m}\cP_1$, and is symmetric, $\hat{R}^{ij}_{kl}=\hat{R}_{ij}^{kl}$. The matrix $\hat{R}$ depends on the parameter $q$ and returns to the undeformed case when $q\to1$ ($\lim\limits_{q\to1}\hat{R}^{ij}_{kl}=\delta^j_k\delta^i_l$). The antisymmetric part def\/ines the commutation relations
\begin{gather*}
\cP_A^{ij}x\otimes x=\left(\cP_A\right)^{ij}_{kl}x^kx^l = 0.
\end{gather*}
We will always use the summation convention. The singlet part def\/ines the metric $C_{ij}=C^{ij}$, by $\left(\cP_1\right)^{ij}_{kl}=\frac{C^{ij}C_{kl}}{C}$ with $C=C^{ij}C_{ij}=(1+q^{2-m})[m/2]_{q^2}$. The metric satisf\/ies the relation
\begin{gather}
\label{CR}
C_{jl}\big(\hat{R}^{\pm 1}\big)^{lk}_{st} = \big(\hat{R}^{\mp 1}\big)^{kl}_{js}C_{lt}
\end{gather}
and is its own inverse, $C_{ij}C^{jk}=\delta^k_i$. The braid matrix also satisf\/ies the relation
\begin{gather}
\label{CR2}
C_{ij}\big(\hat{R}^{-1}\big)^{ij}_{kl} = q^{m-1}C_{kl}.
\end{gather}

The generalized norm squared is then def\/ined as $\bold{x}^2=x^iC_{ij}x^j$. This norm squared is central in the algebra $\cO_q$ and is invariant under the co-action of $O_q(m)$. The explicit expressions for the coaction of $O_q(m)$ or the dually related action of $\cU_q(\mathfrak{so}(m))$ can be found in e.g.~\cite{MR1138053, MR1239953, MR1328733, MR2006509}. In order to obtain a Fourier transform a second set of coordinates is needed, denoted by $\bold{y}$, which is a copy of the $\bold{x}$ coordinates. The commutation relations between the $\bold{x}$ and $\bold{y}$ coordinates are given by $y^ix^j=q^{-1}\hat{R}^{ij}_{kl}x^ky^l$, see~\cite{MR1303081, MR1235979, MR1317458}.

The dif\/ferential calculus on $\mR^m_q$ was developed in~\cite{MR1138053, MR1182031}, the action of the partial derivatives is determined by
\begin{gather}
\label{defpart}
\partial^ix^j = C^{ij}+q\big(\hat{R}^{-1}\big)^{ij}_{kl}x^k\partial^l .
\end{gather}
The Laplace operator on $\mR^m_q$ is given by $\Delta=\partial^iC_{ij} \partial^j$. It is central in the algebra generated by the partial derivatives and is $O_q(m)$-invariant. The commutation relations for the partial derivatives can be expressed using $\cP_A$ or as
\begin{gather}
\label{Rxx}
\hat{R}^{ij}_{kl}\partial^k\partial^l = q\partial^i\partial^j+\frac{1-q^2}{q^{m-1}(1+q^{2-m})} C^{ij}\Delta,
\end{gather}
see \cite{MR1182031}. Formulas \eqref{CR} and \eqref{CR2} yield
\begin{gather}
\label{Laplx}
\Delta x^j=\mu\partial^j+q^2x^j\Delta \qquad \mbox{and} \qquad \partial^j\bold{x}^2=\mu x^j+q^2\bold{x}^2\partial^j\qquad  \mbox{with} \quad\mu=1+q^{2-m}.
\end{gather}
The dilatation operator is given by
\begin{gather}
\label{dil}
\Lambda=1+\big(q^2-1\big)x^iC_{ij}\partial^j+\frac{(q^2-1)^2}{q^{m-2}\mu^2}\bold{x}^2\Delta
\end{gather}
and satisf\/ies $\Lambda x^i=q^2x^i\Lambda $. For $u\in\mR$, $\Lambda^u$ is def\/ined by $\Lambda^u x^i=q^{2u}x^i\Lambda^u$ and $\Lambda^u(1)=1$.

The elements of $\cO_q$ corresponding to f\/inite summations are the polynomials, the correspon\-ding algebra is denoted by $\cP$. The space $\cP_k$ is def\/ined as the space of the polynomials $P$ in $\cP$ which satisfy
\begin{gather}
\label{propdil}
\Lambda (P)=q^{2k}P \qquad \mbox{or} \qquad \frac{\Lambda-1}{q^2-1}P=[k]_{q^2}P.
\end{gather}
For $f$ analytical in the origin, $f(\bold{x}^2)\in\cO_q$ is def\/ined by the Taylor expansion of $f$. Equa\-tion~\eqref{Laplx} leads to
\begin{gather}
\label{partialr}
\partial^jf\big(\bold{x}^2\big) = x^j\mu\partial_{\bold{x}^2}^{q^2}f\big(\bold{x}^2\big)+f\big(q^2\bold{x}^2\big)\partial^j
\end{gather}
for general functions of $\bold{x}^2$ and the $q$-derivative as def\/ined in formula~\eqref{qder}. The relation{\samepage
\begin{gather}
\label{berpx2}
\partial_{\bold{x}^2}^{q^2}f\big(\bold{x}^2\big) = \left[\frac{1}{(1+q)t}\partial_t^qf\big(t^2\big)\right]_{t^2=\bold{x}^2}
\end{gather}
is a useful calculation rule.}



There exists a second dif\/ferential calculus on $\mR^m_q$ of partial derivatives $\overline{\partial^j}$, which is obtained from the unbarred one by replacing $\partial^j$, $q$, $\hat{R}$, $C$ by $\overline{\partial^j}$, $q^{-1}$, $\hat{R}^{-1}$, $C$. In particular, the relation
\begin{gather}
\label{partialrbar}
\overline{\partial^j}f\big(\bold{x}^2\big) = x^jq^{m-2}\mu\partial_{\bold{x}^2}^{q^{-2}}f\big(\bold{x}^2\big)+f\big(q^{-2}\bold{x}^2\big)\overline{\partial^j}
\end{gather}
holds. The algebra generated by the variables and partial derivatives $\partial^j$ is denoted ${\rm Dif\/f}(\mR^m_q)$. The algebra generated by the variables and partial derivatives $\overline{\partial^j}$ is the same algebra. The polynomial null-solutions of $\overline{\Delta}$ are the same as those of $\Delta$. The space of the null-solution of degree $k$ is denoted by $\cS_k$, so $\cS_k=\cP_k\cap \ker \Delta$.


\begin{definition}
The operator $E$ is given by $E=[\frac{m}{2}]_{q^2}+q^mx^iC_{ik}\partial^k$.
\end{definition}

Using the expression for the dilation operator in formula \eqref{dil} the operator $E$ can be expressed as
\begin{gather*}
E= \left[\frac{m}{2}\right]_{q^2}+q^m\left(\frac{\Lambda-1}{q^2-1}-q^{m-2}\big(q^2-1\big)\bold{x}^2\Delta\right).
\end{gather*}
Property \eqref{propdil} then implies that for $S_k\in\cS_k$,
\begin{gather}
\label{Eharm}
ES_k = \left[\frac{m}{2}+k\right]_{q^2}S_k
\end{gather}
holds. Together with $\Delta$ and $\bold{x}^2$, this operator $E$ generates an algebra which is a $q$-deformation of the universal enveloping algebra of $\mathfrak{sl}_2$.
\begin{theorem}
\label{sl2}
The operators $\Delta/\mu$, $\bold{x}^2/\mu$ and $E$ generate the quantum algebra $\cU_q(\mathfrak{sl}_2)$,
\begin{gather}
\label{commLaplxkwad}
\left[\Delta/\mu,\bold{x}^2/\mu\right]_{q^4}=E,\qquad
%\nonumber
\left[E,\bold{x}^2/\mu\right]_{q^2}=[2]_{q^2} \bold{x}^2/\mu,\qquad
%\nonumber
\left[\Delta /\mu,E\right]_{q^2}=[2]_{q^2} \Delta/\mu.
\end{gather}
\end{theorem}

\begin{proof}\sloppy
Combining equations \eqref{Laplx} and \eqref{CR2} yields equation \eqref{commLaplxkwad}. Equation~\eqref{Laplx} implies $x^iC_{ij}\partial^j\bold{x}^2 =\mu \bold{x}^2+q^2\bold{x}^2x^iC_{ij}\partial^j$, which leads to the second relation. The third relation is calculated in the same way.
\end{proof}

\begin{remark}
As the generators of $\cU_q(\mathfrak{sl}_2)$ are $O_q(m)$-invariant, this quantum algebra and quantum group form the Howe dual pair $(O_q(m),\cU_q(\mathfrak{sl}_2))$, or $(\cU_q(\mathfrak{so}(m)),\cU_q(\mathfrak{sl}_2))$.
\end{remark}

The quantum algebra $\cU_q(\mathfrak{sl}_2)$ is equal to the one in~\cite{MR2595267}. In \cite{MR2595267}, $\cU_q(\mathfrak{sl}_2)$ was generated by the standard Euclidean norm squared $r^2$ on $\mR^m$ and a $q$-deformation of the Laplace operator~$\Delta_q$. Since $\Delta_q$ is still $O(m)$-invariant, the Howe dual pair $(O(m),\cU_q(\mathfrak{sl}_2))$ appeared. Because the $O_q(m)$-invariant harmonic operators on quantum Euclidean space in the present paper generate the same quantum algebra we obtain an important connection between these two theories. In particular the $O_q(m)$-invariant Fourier transform developed in the current paper can be used to construct the $O(m)$-invariant $q$-Fourier transform on Euclidean space, as will be done in Section~\ref{qFeucl}.

\begin{lemma}[Fischer decomposition]
\label{Fischer}
The space $\cP$ decomposes into irreducible pieces under the action of $\cU_q(\mathfrak{so}(m))$ as $($see {\rm \cite{MR2006509, MR1328733})}
\begin{gather*}
\cP=\bigoplus_{j=0}^{\infty}\bigoplus_{k=0}^\infty\bold{x}^{2j}\cS_{k}.
\end{gather*}
\end{lemma}


The operator identities in Theorem~\ref{sl2} yield
\begin{gather*}
%\label{evalE}
E\big(\bold{x}^{2l}S_k\big)=\left(\left[\frac{m}{2}+k+l\right]_{q^2}+q^2[l]_{q^2}\right)\bold{x}^{2l}S_k
\end{gather*}
and $\Delta(\bold{x}^{2l}S_k)=\mu^2[l+k+\frac{m}{2}-1]_{q^2}[l]_{q^2}\bold{x}^{2l-2}S_k$. These calculations and the previous results lead to
\begin{theorem}[Howe duality]\sloppy
The decomposition of $\cP$ into irreducible representations of $\cU_q(\mathfrak{so}(m))$ is given in Lemma~{\rm \ref{Fischer}}. Each space $\bigoplus_j \bold{x}^{2j}\cS_k$ is a lowest weight module of $\cU_q(\mathfrak{sl}_2)$ with weight vectors $\bold{x}^{2j}\cS_k$, the lowest weight vector is $\cS_k$ with weight $[m/2+k]_{q^2}$. The Fischer decomposition of~$\cP$ therefore is a multiplicity free irreducible direct sum decomposition under the joint action of $\cU_q(\mathfrak{sl}_2)\times\cU_q(\mathfrak{so}(m))$.
\end{theorem}

The antilinear involutive antihomomorphism $\ast$ on ${\rm Dif\/f}(\mR^m_q)$ is def\/ined by $(AB)^\ast=B^\ast A^\ast$, $(x^j)^\ast=x^kC_{kj}$,  $(\partial^j)^\ast=-q^{-m}\overline{\partial^k}C_{kj}$ and $\lambda^\ast=\overline{\lambda}$ with $\lambda\in\mC$ and $\overline{\cdot}$ complex conjugation. This yields
\begin{gather*}
\big(\bold{x}^2\big)^\ast=\bold{x}^2 \qquad \mbox{and}\qquad \Delta^\ast=q^{-2m}\overline{\Delta}.
\end{gather*}

The harmonic oscillator on quantum Euclidean space was studied in \cite{MR1097174, MR1288667, MR1239953}. The two Hamiltonians (with an unimportant dif\/ferent normalization compared to~\cite{MR1239953}) are given by
\begin{gather}
\label{HOs}
h=\frac{1}{2}\left(-\Delta+\bold{x}^2\right),\qquad h^\ast=\frac{1}{2}\left(-{\Delta}^\ast+\bold{x}^2\right).
\end{gather}
Both operators have the same eigenvalues.


\subsection{Integration and Fourier theory on quantum spaces}
In \cite{MR1303081} a method was prescribed to generalize $q$-integration to higher dimensions in the context of quantum spaces. Gaussian-induced integration for general $\hat{R}$-matrices is def\/ined assuming there is a matrix $\eta\in\mR^{m\times m}$ and a solution $g_{\eta}\in\cO$ of the equation
\begin{gather}
\label{Gaussian}
-\eta^i_{ j}\partial^jg_{\eta}=x^i g_\eta.
\end{gather}
Integration $\int$ on the space $\cP g_{\eta}$, with $\cP$ the polynomials on the quantum space is then uniquely def\/ined by demanding $\int\circ \partial^i =0$, $i=1,\dots,m$. For $f\in\cP$ the integral $\int f g_{\eta}$ is of the form
\begin{gather}
\label{Gaussian2}
\int fg_{\eta}=Z[f]I(g_\eta),
\end{gather}
with $I(g_\eta)=\int g_\eta$ and $Z$ a functional on~$\cP\subset\cO$. Superspace with purely bosonic and fermionic coordinates can be seen as a limit of braided spaces, typically for $q\to-1$, see e.g.~\cite{MR1413908}. From this point of view it is interesting to note that the Berezin integral can also be constructed in this setting. In \cite{DBS5, MR2539324} this led to integration over the supersphere and a new interpretation of the Berezin integral.

An explicit example of this construction was already def\/ined on quantum Euclidean space, see~\cite{MR1239953}. In~\cite{MR1408120} it was shown that this integration can be def\/ined and generalized using integration over the quantum Euclidean sphere. In Section~\ref{DefFourier} we will show how this approach follows from harmonic analysis on quantum Euclidean space.


In \cite{MR1303081} a general procedure to construct a Fourier transform on quantum spaces was developed. First the appropriate Gaussian-induced integration $\int$ should be constructed and the exponential or Fourier kernel (see \cite{MR1235979, MR2099760}) calculated. The Fourier transform on a braided-Hopf algebra $B$ with left dual Hopf algebra $B^\ast$  is a map $\cF:B\to B^\ast$. The co-ordinates for $B$ are denoted by $\bold{x}$ and for $B^\ast$ by $\bold{y}$, the Fourier transforms are given by
\begin{gather*}
\cF[f(\bold{x})](\bold{y})=\int_x f(\bold{x})\exp_{\hat{R}}(\bold{x}|\bold{y}), \qquad \cF^\ast[f(\bold{y}](\bold{x})=\int^\ast_y f(\bold{y})\exp_{\hat{R}}(\bold{x}|\bold{y}).
\end{gather*}
These Fourier transforms are each others inverse, $\cF^\ast\cF=\mbox{Vol}\, S$, with $S$ the antipode on $B$. As an explicit example we consider the braided line $B=\mC[x]$, with braided-Hopf algebra structure as introduced in \cite{MR1182915} given by
\begin{gather*}
\Delta x^k=\sum_{j=0}^k\frac{[k]_q!}{[k-j]_q![j]_q!}x^j\otimes x^{k-j},\qquad Sx^k=(-1)^kq^{\frac{k(k-1)}{2}}x^k,\qquad\epsilon x^k=\delta_{k0}.
\end{gather*}
The dually-paired Hopf algebra $B^\ast$ is the same Hopf algebra with variable~$y$. The relation $xy=qyx$ holds. The exponential is $\exp(x|y)=\sum\limits_{k=0}^\infty\frac{x^ky^k}{[k]_q!}$ and satisf\/ies $\partial_x^q\exp(x|y)=\exp(x|y)y$. The Fourier transforms take the form
\begin{gather}
\label{Fourline}
\cF[f](y)=\int_{-\gamma\cdot\infty}^{\gamma\cdot\infty}d_qxf(x)\exp(x|y) \qquad \mbox{and}\qquad \cF^\ast[f](x)=\int_{-\delta\cdot\infty}^{\delta\cdot\infty}d_qyf(y)\exp(x|y).
\end{gather}
The theory of \cite{MR1303081} then implies $\cF^\ast\cF[f](x)=Sf(x)\mbox{Vol}_{\gamma,\delta}$. We can rewrite this in a way that will be more closely related to our approach of the Fourier transform on $\mR^m_q$. Def\/ine $g(y)=\frac{1}{{\rm Vol}_{\gamma,\delta}}\cF[f](y)$, then the def\/inition of the antipode implies
\begin{gather*}
f(-x) = \int_{-\delta\cdot\infty}^{\delta\cdot\infty}d_qyg(y)\sum_{k=0}^\infty\frac{q^{\frac{k(k+1)}{2}}y^kx^k}{[k]_q!}.
\end{gather*}
In this equation and in the f\/irst equation of~\eqref{Fourline}, there are no coordinates which have to be switched before integration. This implies that we can assume $x$ and $y$ commute and write the equations above as
\begin{gather*}
g(y)=\frac{1}{{\rm Vol}_{\gamma,\delta}}\int_{-\gamma\cdot\infty}^{\gamma\cdot\infty}d_qx\,f(x)\,e_q(xy)\qquad \mbox{and}\qquad f(x)=\int_{-\delta\cdot\infty}^{\delta\cdot\infty}d_qy\,g(y)\,E_q(-qyx).
\end{gather*}



In \cite{MR1448685} a closely related analytical approach was given to the one dimensional Fourier transform above. Consider real commuting variables~$x$ and~$y$. Using the orthogonality relations and the generating function of the Hermite polynomials in equation~\eqref{Hermite}, it is possible to prove
\begin{gather}
\frac{1}{2\Gamma_{q^2}(\frac{1}{2})}\int_{-\frac{1}{\sqrt{1-q}}}^{\frac{1}{\sqrt{1-q}}}d_qx\, \left[H_k^q\left(\frac{x}{\sqrt{1+q}}\right)E_{q^2}\left(-\frac{x^2}{1+q}\right)\right]e_q(-ixy)\nonumber\\
\qquad{}
=\frac{(q+1)^{\frac{k-1}{2}}}{i^k}q^{\frac{1}{2}(k+1)(k+2)}y^ke_{q^2}\left(-\frac{q^2y^2}{1+q}\right),\label{1dFour}
\end{gather}
which is equivalent to~\cite[equation~(8.7)]{MR1448685}. For every $\delta\in\mR^+$,
\begin{gather}
\frac{1}{C_\delta}\int_{-\delta\cdot\infty}^{\delta\cdot\infty}d_qy\, \left[y^ke_{q^2}\left(-\frac{q^2y^2}{1+q}\right)\right]E_q(iqyx)\nonumber\\
\qquad =\frac{i^k}{q^{\frac{1}{2}(k+1)(k+2)}(1+q)^{\frac{k-1}{2}}}H_k^q\left(\frac{x}{\sqrt{1+q}}\right)E_{q^2}\left(-\frac{x^2}{1+q}\right)\label{1dFourb}
\end{gather}
holds for $C_\delta$ some constant depending on $\delta$, see~\cite[equation~(8.21)]{MR1448685}. So the two Fourier transforms as def\/ined in equations~\eqref{1dFour} and~\eqref{1dFourb} can be regarded as each others inverse, which was to be expected from the theory of~\cite{MR1303081}. There is however one dif\/ference between the explicit Fourier transform in~\cite{MR1448685} and the abstract theory in~\cite{MR1303081}. While the inverse Fourier transform remains unchanged, the integration for the Fourier transform is limited to a f\/inite interval. This will be explained in the subsequent Lemma~\ref{intexpinf}. However, using property~\eqref{intinfeind} the integral can be replaced by $\int_{-\frac{1}{\sqrt{1-q}}\cdot\infty}^{\frac{1}{\sqrt{1-q}}\cdot\infty}$. So the analytical approach of~\cite{MR1448685} recovers the theory from~\cite{MR1303081} with an imposed limitation on~$\gamma$. For other~$\gamma$, the theory from~\cite{MR1303081} would still hold, but the constant ${\rm Vol}_{\gamma,\delta}$ is inf\/inite.



\section[The $q$-Hankel transforms]{The $\boldsymbol{q}$-Hankel transforms}
\label{Hanksect}

In this section we def\/ine two $q$-Hankel transforms using the f\/irst and second $q$-Bessel function. These transforms will act as each others inverse. This is a generalization of the result in formulas~\eqref{1dFour} and~\eqref{1dFourb}. By evaluating the Fourier transform on the appropriate functions in Section~\ref{boch} the f\/irst of these $q$-Hankel transforms will appear.

We will calculate the $q$-Hankel transforms of the $q$-Laguerre and little $q$-Laguerre polynomials weighted with a $q$-Gaussian. The undeformed Fourier--Gauss transform of these polynomials was already studied in \cite{MR1479780}. There also exists a third $q$-Bessel function besides the ones in equations~\eqref{Bessel1} and~\eqref{Bessel2}. We will not explicitly need the third type, but it is interesting to note that in~\cite{MR1069750, MR2445269} the corresponding $q$-Hankel transforms were constructed. These Hankel transforms could also be used to def\/ine an $O_q(m)$-invariant Fourier transform on $\mR^m_q$. This would have the advantage that the Fourier transform is its own inverse. That Fourier transform would however not behave well with respect to the derivatives on $\mR^m_q$. This is already the case for the Fourier transform on the braided line, as is proven in~\cite{MR1069750}. The braided line corresponds to $\mR^m_q$ for $m=1$.


In anticipation of the connection with the Fourier transform on quantum Euclidean space we will scale the $q$-Hankel transforms in the following def\/inition with~$\mu$, see equation~\eqref{Laplx}, although at this stage any constant could be used. The reason for the appearance of unf\/ixed constants~$\beta$,~$\gamma$ will become apparent in Sections~\ref{boch}, \ref{behave} and \ref{FourHO}.
\begin{definition}
\label{defHankel}
For $\beta,\gamma\in\mR^+$ and $\nu\ge-\frac{1}{2}$, the $q$-Hankel transforms are given by
\begin{gather*}
\overline{\cH}_{\nu}^{q,\beta}[f(r)](t) = \frac{1+q}{\mu}\int_0^{\sqrt{\frac{\mu}{(1-q^2)\beta}}}d_qr\frac{J_\nu^{(1)}\big(\frac{1+q}{\mu}rt|q^2\big)}{(rt)^{\nu}}r^{2\nu+1}[f(r)]
\end{gather*}
and
\begin{gather*}
\cH^{q,\gamma}_{\nu}[f(t)](r) = \frac{1+q}{\mu}\int_0^{\gamma\cdot\infty}d_qt\frac{J_\nu^{(2)}\big(q\frac{1+q}{\mu}rt|q^2\big)}{(rt)^{\nu}}t^{2\nu+1}[f(t)].
\end{gather*}
\end{definition}

In this def\/inition it is not specif\/ied on which function spaces the $q$-Hankel transforms act. At the moment we def\/ine them on functions for which the expression exists.

In order to connect the Fourier transform on $\mR^m_q$ with these $q$-Hankel transforms we def\/ine the following transformations,
\begin{gather*}
\overline{\cF_{\nu}}^{q,\beta}[\psi]\big(t^2\big)=\overline{\cH}_{\nu}^{q,\beta}[\psi\circ\Upsilon](t)
\qquad \mbox{and} \qquad \cF_{\nu}^{q,\gamma}[\psi]\big(r^2\big)=\cH_{\nu}^{q,\gamma}[\psi\circ\Upsilon](r),
\end{gather*}
with $\Upsilon(u)=u^2$.





In order to prove the properties of the $q$-Hankel transforms we will need some identities of the $q$-Bessel functions in equations~\eqref{Bessel1} and~\eqref{Bessel2}. These are summarized in the following lemma.

\begin{lemma}
\label{lemBessel1}
%\label{lemBessel2}
The first and second $q$-Bessel functions satisfy
\begin{alignat*}{3}
& (i)& &\partial_u^q\frac{J_{\nu}^{(1)}(u|q^2)}{u^\nu}=-u\frac{J_{\nu+1}^{(1)}(u|q^2)}{u^{\nu+1}},\qquad \partial_u^{q^{-1}}\frac{J_{\nu}^{(2)}(qu|q^2)}{u^\nu}=-qu\frac{J_{\nu+1}^{(2)}(qu|q^2)}{u^{\nu+1}}, &\\
& (ii)& &\frac{J_{\nu+1}^{(1)}(u|q^2)+J_{\nu-1}^{(1)}(u|q^2)}{u^{\nu-1}}=[2\nu]_q\frac{J_{\nu}^{(1)}(qu)}{(qu)^\nu},& \\ &&& \frac{J_{\nu+1}^{(2)}(qu|q^2)+J_{\nu-1}^{(2)}(qu|q^2)}{u^{\nu-1}}=[2\nu]_q\frac{J_{\nu}^{(2)}(u)}{(qu)^\nu},&\\
& (iii) \ \ & &\partial_u^q J_{\nu}^{(1)}(u|q^2)u^\nu=J_{\nu-1}^{(1)}(u|q^2)u^\nu\qquad
 \mbox{and}\qquad \partial_u^q J_{\nu}^{(2)}(u|q^2)u^\nu=q^\nu J_{\nu-1}^{(2)}(qu|q^2)u^\nu.&
\end{alignat*}
\end{lemma}

\begin{proof}
The right-hand side of the second property (for the f\/irst $q$-Bessel function) can be calculated using $[\nu]_{q^2}q^{2i}=[\nu+i]_{q^2}-[i]_{q^2}$,
\begin{gather*}
[\nu]_{q^2}(1+q)\frac{J_\nu^{(1)}(qu)}{(qu)^\nu} \\
\qquad{}
=\frac{1}{(1+q)^{\nu-1}}\sum_{i=0}^\infty\left(\frac{(-1)^i}{[i]_{q^2}!\Gamma(i+\nu)}\left(\frac{u}{1+q}\right)^{2i}
-\frac{(-1)^i}{[i-1]_{q^2}!\Gamma(i+\nu+1)}\left(\frac{u}{1+q}\right)^{2i}\right)\\
\qquad{} =\frac{J_{\nu-1}^{(1)}(u|q^2)+J_{\nu+1}^{(1)}(u|q^2)}{u^{\nu-1}}.
\end{gather*}
The f\/irst and the third property follow from a direct calculation. The left-hand sides of the properties can also be obtained from \cite[Exercise~1.25]{MR1052153}.
\end{proof}


\begin{corollary}
\label{lemBesselextra}
The second $q$-Bessel functions satisfy the following relation:
\begin{gather*}
\partial_u^q u^{\nu+1}J_{\nu-1}^{(2)}\big(u|q^2\big) = [2\nu]_qu^\nu J_{\nu-1}^{(2)}\big(u|q^2\big)-q^{\nu+1}u^{\nu+1}J_{\nu}^{(2)}\big(qu|q^2\big).
\end{gather*}
\end{corollary}
\begin{proof}
This is a direct consequence of the second formula in Lemma~\ref{lemBessel1}$(i)$.
\end{proof}




Combining generating function \eqref{BesselLag1} and orthogonality relation \eqref{orthLag1} yields
\begin{gather}
\int_0^{\frac{1}{\sqrt{1-q}}}d_qr \frac{J_\nu^{(1)}(rt|q^2)}{(rt)^{\nu}}r^{2\nu+1}\left[\cL_j^{(\nu)}\left(\frac{r^2}{1+q}|q^2\right)E_{q^2}\left(\frac{-r^2}{1+q}\right)\right]
\nonumber\\
\qquad {} =\frac{q^{2(j+1)(j+\nu+1)}}{[j]_{q^2}!}\frac{t^{2j}}{(1+q)^j}e_{q^2}\left(-\frac{q^2t^2}{1+q}\right).\label{preHank1}
\end{gather}
Generating function \eqref{BesselLag2} and orthogonality relation \eqref{orthLag2} lead to
\begin{gather} \int_0^{\gamma\cdot\infty}d_qt\frac{J_\nu^{(2)}(qrt|q^2)}{(rt)^{\nu}}t^{2\nu+1}\left[\cL_j^{(\nu)}
\left(\frac{q^2t^2}{1+q}|q^{-2}\right)e_{q^2}\left(\frac{-q^2t^2}{1+q}\right)\right]\nonumber\\
\qquad {}= d\left(\frac{\gamma}{\sqrt{1+q}},\nu\right)\frac{q^{-(j+1)(j+2\nu+2)}}{[j]_{q^2}!}\frac{r^{2j}}{(1+q)^j}E_{q^2}
\left(-\frac{r^2}{1+q}\right).\label{preHank2}
\end{gather}

The following expansion of a monomial in terms of the $q$-Laguerre polynomials is a direct consequence of equation \eqref{BesselLag2},
\begin{gather*}
\frac{t^{2j}}{(1+q)^j[j]_{q^2}!}=\sum_{i=0}^j\frac{(-1)^i(q^{2i+2\nu+2};q^2)_{(j-i)}}{[j-i]_{q^2}!(1-q^2)^{j-i}}
\frac{q^{(j-i)(j-i+1)+(i+1)(i+2\nu+2)}}{q^{2(j+1)(j+\nu+1)}}\cL_i^{(\nu)}\left(\frac{q^2t^2}{1+q}|q^{-2}\right).
\end{gather*}
Applying this yields
\begin{gather*}
\int_0^{\gamma\cdot\infty}d_qt\frac{J_\nu^{(2)}(qrt|q^2)}{(rt)^{\nu}}t^{2\nu+1}
\left[\frac{t^{2j}}{(1+q)^j[j]_{q^2}!}e_{q^2}\left(\frac{-q^2t^2}{1+q}\right)\right]\\
\quad{}=\sum_{i=0}^j\frac{(-1)^i(q^{2i+2\nu+2};q^2)_{(j-i)}}{[j-i]_{q^2}!(1-q^2)^{j-i}}\frac{q^{(j-i)(j-i+1)}}{q^{2(j+1)(j+\nu+1)}}
d\left(\frac{\gamma}{\sqrt{1+q}},\nu\right)\frac{1}{[i]_{q^2}!}\frac{r^{2i}}{(1+q)^i}E_{q^2}\left(-\frac{r^2}{1+q}\right)\\
\quad{}=d\left(\frac{\gamma}{\sqrt{1+q}},\nu\right) q^{-2(j+1)(j+\nu+1)}\cL^{(\nu)}_j\left(\frac{r^2}{1+q}|q^2\right)E_{q^2}\left(-\frac{r^2}{1+q}\right).
\end{gather*}
These calculations imply the following relations for the $q$-Hankel transforms in Def\/inition~\ref{defHankel}:
\begin{gather*}
\overline{\cH}_{\nu}^{q,\beta}\left[\cL_j^{(\nu)}\left(\beta\frac{r^2}{\mu}|q^2\right)E_{q^2}\left(-\beta\frac{r^2}{\mu}\right)\right](t)=
\frac{1}{\beta^{\nu+1+j}}C_jt^{2j}e_{q^2}\left(-\frac{q^2t^2}{\mu\beta}\right)\qquad\mbox{for} \ \ \beta\in\mR^+,\\
\cH^{q,\gamma}_{\nu}\left[C_jt^{2j}e_{q^2}\left(-\alpha\frac{q^2t^2}{\mu}\right)\right](r)= \frac{d\big(\frac{\sqrt{\alpha}\gamma}{\sqrt{\mu}},\nu\big) }{\alpha^{\nu+1+j}}\cL^{(\nu)}_j\left(\frac{r^2}{\alpha\mu}|q^2\right)E_{q^2}\left(-\frac{r^2}{\alpha\mu}\right)\!\qquad\mbox{for} \ \ \alpha,\gamma\in\mR^+,
\end{gather*}
with $C_j=\frac{q^{2(j+1)(j+\nu+1)}}{[j]_{q^2}!\mu^j}$.


By considering the case $\beta=1/\alpha$ we obtain
\begin{theorem}
\label{Hankelinv}
For each $\alpha,\gamma\in\mR^+$, the inverse of the $q$-Hankel transform $\cH_\nu^{q,\gamma}$ acting on $\mR[t^2]\otimes e_{q^2}(-\alpha\frac{q^2t^2}{\mu})$,
\begin{gather*}
\cH_\nu^{q,\gamma}: \ \ \mR\big[t^2\big]\otimes e_{q^2}\left(-\alpha\frac{q^2t^2}{\mu}\right) \to\mR\big[r^2\big]\otimes E_{q^2}\left(-\frac{r^2}{\alpha\mu}\right)
\end{gather*}
is given by
\begin{gather*}
\frac{1}{d\big(\frac{\sqrt\alpha \gamma}{\sqrt\mu},\nu\big)}\overline{\cH}^{q,1/\alpha}_{\nu}: \ \ \mR\big[r^2\big]\otimes E_{q^2}\left(-\frac{r^2}{\alpha\mu}\right)\to \mR\big[t^2\big]\otimes e_{q^2}\left(-\alpha\frac{q^2t^2}{\mu}\right).
\end{gather*}
\end{theorem}
This theorem for $\nu=-\frac{1}{2}$ and $\nu=\frac{1}{2}$ is identical to the results in equations \eqref{1dFour} and \eqref{1dFourb}.


In order to prove the behavior of the Fourier transform on $\mR^m_q$ we need the following properties of the $q$-Hankel transforms. The exact function spaces on which they act is again not specif\/ied, the properties hold if all the terms are well-def\/ined. In particular these lemmata hold for the function spaces in Theorem~\ref{Hankelinv}.

